\section{Domain-Specific Notation}
\label{app:domain-specific-notation}

\begin{definition}[SI unit convention]
\label{def:domain-si-unit-convention}
Unless otherwise stated, physical continuum quantities use SI units: length in
meters, time in seconds, density in $\mathrm{kg}/\mathrm{m}^3$, pressure in Pa,
dynamic viscosity in $\mathrm{Pa}\cdot\mathrm{s}$, and kinematic viscosity in
$\mathrm{m}^2/\mathrm{s}$.
\end{definition}

\begin{convention}[Centimeter--second numerical-record convention]
\label{conv:cm-second-numerical-record}
The Julia solver and the generated report tables use the centimeter--gram--second
numerical-record convention unless a table states otherwise. Lengths are in
$\mathrm{cm}$, time in $\mathrm{s}$, density in $\mathrm{g}/\mathrm{cm}^3$,
velocity in $\mathrm{cm}/\mathrm{s}$, pressure and stress in
$\mathrm{dyn}/\mathrm{cm}^2$, dynamic viscosity in
$\mathrm{g}/(\mathrm{cm}\,\mathrm{s})$, and kinematic viscosity in
$\mathrm{cm}^2/\mathrm{s}$.

The implemented 1D solver stores a radius-squared coordinate $a=R^2$ and a
solver flow $q$. Physical circular area and physical flow are
\begin{flalign*}
    \quad A_{\mathrm{phys}}=\pi a,
    \qquad Q_{\mathrm{phys}}=\pi q,
    \qquad \bar u=\frac{q}{a}=\frac{Q_{\mathrm{phys}}}{A_{\mathrm{phys}}}. &&
\end{flalign*}
Thus legacy CSV headers such as \texttt{A\_cm2} and \texttt{Q\_cm3\_s} are read
in this report as the solver-coordinate fields $a$ and $q$ unless a physical
quantity is explicitly labelled \texttt{Aphys\_cm2} or \texttt{Qphys\_cm3\_s}.
A Pa-valued command-line pressure is converted before storage by
$1\,\mathrm{Pa}=10\,\mathrm{dyn}/\mathrm{cm}^2$.
\end{convention}

\begin{convention}[Pressure-output convention]
\label{conv:appendix-pressure-reference-ratio}
Pressure-drop and pressure-ratio symbols are retained only as reduced-output
vocabulary. A pressure-drop trace requires declared proximal and distal scalar
pressure observations. A pressure-ratio output additionally requires a declared
observation tuple
\begin{flalign*}
    \quad
    \pi_{\mathrm{PR}}
    \defined
    (\mathcal O_{\mathrm{prox}},
    \mathcal O_{\mathrm{dist}},
    \mathcal T_{\mathrm{HF}},
    p_{\mathrm{ref}}), &&
\end{flalign*}
with a finite averaging window and a positive denominator admissibility bound.
The scalar observations in this tuple may be section averages, reconstructed
wall-law representatives, or other declared pressure observations; they are
pointwise pressures only when the corresponding observation map is a
point-evaluation or point-reconstruction map.
For observed gauge-pressure traces $P_{\mathrm{prox}}$ and $P_{\mathrm{dist}}$,
\begin{flalign*}
    \quad
    \Delta p(t) &\defined P_{\mathrm{prox}}(t)-P_{\mathrm{dist}}(t), &&\\
    \quad
    \mathrm{PR}_{1D}(\pi_{\mathrm{PR}}) &\defined
    \frac{\langle p_{\mathrm{ref}}+P_{\mathrm{dist}}\rangle_{\mathcal T_{\mathrm{HF}}}}
         {\langle p_{\mathrm{ref}}+P_{\mathrm{prox}}\rangle_{\mathcal T_{\mathrm{HF}}}}. &&
\end{flalign*}
The subscript $\mathrm{HF}$ is only a declared averaging-window label in this
report, and no pressure-ratio result is reported in the diagnostic velocity
comparison. The shorthand is not a clinical measurement model
\cite[Ch.~5, Sec.~5.5, pp.~72--73]{EscanedDavies2017PhysiologicalAssessment}
\cite[p.~1703]{PijlsEtAl1996FFRFunctionalSeverity}.
\end{convention}

\begin{convention}[Shear-output convention]
\label{conv:appendix-shear-diagnostic}
Shear terminology is model-tier dependent. A resolved wall-shear-stress output
requires a resolved stress tensor, wall normal, wall location, output operator,
and, for a signed component, a tangent direction
\cite[Abstract; Sec.~2, pp.~8--10]{JohnEtAl2017WallShearStressVectorForm}. A
reduced 1D shear proxy requires its closure, profile, location, and sign or
window convention. Reduced shear proxies are closure diagnostics, not resolved
traction outputs.
\end{convention}

\noindent\begin{minipage}[t]{\textwidth}
\begin{compactacronymtables}
\setlength{\parskip}{0pt}
\begin{tabular}{@{}>{\raggedright\arraybackslash}p{0.24\linewidth}
        >{\raggedright\arraybackslash}p{0.56\linewidth}
        >{\raggedright\arraybackslash}p{0.14\linewidth}@{}}
        \toprule
        \textbf{Symbol} & \textbf{Meaning} & \textbf{Units} \\
        \midrule
        $\vu(t,\vx)$, $p(t,\vx)$ & continuum velocity and pressure field values when a continuum tier is being referenced & length/time; pressure \\
        $\vT(t,\vx)$, $\tau_{ij}$ & continuum Cauchy stress tensor and stress components before reduction or observation & pressure \\
        $R$, $d_v=2R$ & vessel radius and diameter & length \\
        $R_0(z)$ & reference stenosed radius used by the geometry and wall source & length \\
        $R_{\max}$ & healthy reference-radius parameter used by the implemented $R_{\max}$-normalized wall law & length \\
        $a=R^2$ & solver radius-squared area coordinate stored as \texttt{A} in legacy output & length$^2$ \\
        $q$ & solver flow coordinate stored as \texttt{Q} in legacy output & length$^3$/time \\
        $A_{\mathrm{phys}}=\pi a$ & physical circular cross-sectional area & length$^2$ \\
        $Q_{\mathrm{phys}}=\pi q$ & physical circular-area flow & length$^3$/time \\
        $\bar u=q/a$ & section-mean velocity used by the 1D--3D observation comparison & length/time \\
        $\alpha_{\mathcal P}$ & momentum-flux correction for the declared velocity profile; the parabolic baseline uses $4/3$ & -- \\
        $g_{\mathcal P}$ & profile shear-rate factor used in the reduced wall-friction and viscosity proxy & -- \\
        $\mathcal R$ & rheology closure family used by the reduced closure & mixed \\
        $\mathcal W$ & wall closure family; the selected report model uses the Canic--Koiter law & mixed \\
        $K=Eh/(1-\sigma^2)$ & wall stiffness coefficient used by the implemented wall law & force/length \\
        $\beta(z)$ & wall-law stiffness parameter; for the selected evolution law $\beta=\sqrt\pi K R_0^2/R_{\max}^2$ & pressure\,$\cdot$\,length \\
        $\Psi$ & elastic flux potential satisfying $\pd_A\Psi=(A/\rho)\pd_Ap_{\mathcal W,\mathrm g}^{\mathrm{evol}}$ in the physical-area notation & length$^4$/time$^2$ \\
        $\mathbf U=(a,q)^\top$ & implemented finite-volume state in solver coordinates & mixed \\
        $\widehat{\mathbf F}_{i+1/2}$ & numerical flux at a finite-volume interface & mixed \\
        $\widehat S_i$ & cell-centered source term used by the implemented operator & mixed \\
        $D_z^h$ & recorded finite-difference operator on cell-centered quantities & 1/length \\
        $\dot\gamma_{1D}=g_{\mathcal P}|q|/(a\sqrt a)$ & reduced shear-rate proxy used for generalized-Newtonian computed closure-family cases & 1/time \\
        $\nu_{\mathrm{eff}}^{\mathcal R}$ & density-scaled effective kinematic viscosity used in reduced source terms & length$^2$/time \\
        $c$ or $c_{\mathrm{wave}}$ & 1D pulse-wave speed for the declared wall law & length/time \\
        $p_{\mathcal W,\mathrm g}^{\mathrm{evol}}$ & reduced gauge pressure produced by the selected evolution wall law & pressure \\
        $p_{\mathrm{loc},\mathrm g}^{\mathrm{diag}}$ & legacy local-denominator pressure diagnostic from the postprocessing routine & pressure \\
        $\tau_{w,z}^{1D}$ & signed 1D axial shear proxy after declaring closure and output convention & pressure \\
        $\mathcal G_{1D,h}^{r_h,e}$ & computed fixed-grid 1D map indexed by numerical and closure records & mixed \\
        \bottomrule
\end{tabular}
\end{compactacronymtables}
\end{minipage}

\begin{remark}[Viscosity notation]
\label{rmk:viscosity-notation}
This report writes $\eta$ for dynamic viscosity and
\begin{flalign*}
    \quad \nu:=\frac{\eta}{\rho} &&
\end{flalign*}
for kinematic viscosity. Generalized-Newtonian closure families use a stress-level
effective dynamic viscosity $\eta_{\mathrm{eff}}(\dot\gamma)$ and the
corresponding density-scaled coefficient
$\nu_{\mathrm{eff}}^{\mathcal R}=\eta_{\mathrm{eff}}^{\mathcal R}/\rho$ in
reduced source terms. This convention avoids overloading $\nu$ with the dynamic
coefficient used in continuum stress laws
\cite[Sec.~3.1, p.~35; Sec.~3.6, p.~215]{KopplHelmig2023DimensionReducedModeling}
\cite[Sec.~1.1, pp.~6--7; Eqs.~(1.2)--(1.4); Fig.~4.4, p.~18]{JohnEtAl2017WallShearStressVectorForm}.
\end{remark}
